\documentclass[11pt]{article}
\usepackage{fontspec}
\usepackage[french,english]{babel}
\usepackage{amsmath}
\usepackage{amssymb}
\usepackage{amsthm}
\usepackage{mathrsfs}
\usepackage{tikz-cd}
\usepackage{enumitem}
\usetikzlibrary{babel}

\newcommand{\printedpage}[1]{\marginpar{\footnotesize p.~#1}}
\newcommand{\msp}[1]{\marginpar{\footnotesize [#1]}}
\newcommand{\uGamma}{\underline{\Gamma}}
\newcommand{\Hom}{\operatorname{Hom}}
\newcommand{\Ob}{\mathrm{Ob}\,}
\newcommand{\card}{\mathrm{card}\,}

\begin{document}

\printedpage{94}
% head: A. Grothendieck : Dérivateurs, Ch. XVIII, Édition des Universités de Montpellier II et Paris VII
thus
\[
\begin{aligned}
\Hom_{\mathcal{M}}(x,X) &\simeq \varinjlim_{I'} \Hom(x,\xi_{i'})\\
&\simeq \varinjlim_{I'} \varprojlim_{I} \Hom(x_i,\xi_{i'})\\
&\overset{\text{SGA 4 I 9.8}}{\simeq} \varprojlim_{I} \underbrace{\varinjlim_{I'} \Hom(x_i,\xi_{i'})}_{\simeq \Hom(x_i,X)}\\
&\simeq \varprojlim_{I} \Hom(x_i,X)\ ,
\end{aligned}
\]
which proves that $x$ is also a $\varinjlim$ of the $x_i$ in the whole of $\mathcal{M}$.

\msp{250}
Thus, we are reduced to proving 4.10.3. Applying 4.9.6 2°, we are brought back to proving, for the functor under consideration, conditions F 1 and F $3_\pi$ (page 225).

In the arguments that follow, we are going to use conditions a), b) of 4.10.3, but not, for the moment, condition c.

\textbf{Lemma 4.10.5.} \textit{Under the conditions of 4.10.3 a) b) (excluding c), for (4.10.3.1) to be cofinal and $\Gamma(\mathcal{E}/J)$ large relative to $\pi$, it is necessary and sufficient that the functor under consideration satisfy condition F 1, recalled below, and its variant F$'$ 1 stated below:}
\begin{itemize}[leftmargin=3em]
\item[F\,1)] {\itshape For every system $(X(j))_{j\in\Ob J}$ of elements $X(j)\in\Ob\mathcal{E}_j$, exists a \textup{section} $G$ of $\mathcal{E}$ over $J$, and a system of arrows
\[
X(j) \longrightarrow G(j) \qquad (j\in ObJ)\ . \tag*{\textup{(4.10.5.1)}}
\]}
\item[F$'$1)] {\itshape Under the conditions of the preceding \textup{F\,1}, suppose in addition a section $F$ of $\mathcal{E}$ over $J$ \msp{251} and a system of arrows
\[
F(j) \longrightarrow X(j) \qquad (j\in\Ob J)\ ,
\]
then $G$ and the system of arrows (4.10.5.1) can be chosen in such a way that the composite arrows
\[
F(j) \longrightarrow X(j) \longrightarrow G(j)
\]
define a homomorphism of sections
\[
F \longrightarrow G
\]
(in the category $\uGamma(\mathcal{E}/J)$).}
\end{itemize}
(\footnote[106]{\textbf{N.B.} Conditions F 1 and F$\,'$\,1 are equivalent to saying that the functor under consideration satisfies conditions F1 and F3, seeing also (by virtue of 4.9.6.3, 1°) that $\uGamma(\mathcal{E}/J)$ is \textit{filtering} without more, and the faithful functor (4.10.3.1) cofinal. This already implies that $\uGamma(\mathcal{E}/J)$ is large relative to $\pi$.}).

\textsc{Proof}. We have to see that F 1 \& F$'$ 1 is equivalent to F 1 \& F $3_\pi$. The implication $\Longleftarrow$ is $\pm$ immediate, it is $\Longrightarrow$ that we shall need, so
\[
\text{F 1 \& F}'\text{ 1} \quad\Longrightarrow\quad \text{F } 3_\pi\ .
\]
% foot: Version 17 décembre 2025

\printedpage{95}
% head: A. Grothendieck : Dérivateurs, Ch. XVIII, Édition des Universités de Montpellier II et Paris VII
For this, let us restate F $3_\pi$ in the present context. We are given a family $(F_\alpha)_{\alpha\in A}$, $\card A\leqslant\pi$, of objects of $\uGamma(\mathcal{E}/J)$, an object $F$ of the same category, and for every $j$ in $\Ob J$, a system of arrows
\[
F_\alpha(j) \xrightarrow{u_\alpha(j)} F(j)\ ,
\]
but without supposing that for $\alpha$ fixed, $j$ variable

\msp{252}
this defines a homomorphism of sections $F_\alpha\longrightarrow F$. We seek a section $G$ of $\mathcal{E}$ over $J$, and a homomorphism of sections $F\longrightarrow G$, such that for every $\alpha$, the system of composites
\[
F_\alpha(j) \longrightarrow F(j) \longrightarrow G(j)
\]
defines a homomorphism of sections $F_\alpha\longrightarrow G_\alpha$. So let $\lambda : j\longrightarrow j'$ be an arrow of $J$, let us consider the squares (\textit{noncommutative}) % [sic] F_\alpha -> G_\alpha -> F_\alpha -> G, the section sought is G (cf. p. 95 below: F_\alpha -> G)
\[
\begin{tikzcd}
F_\alpha(j) \arrow[r,"u_\alpha(j)"] \arrow[d,"F_\alpha(\lambda)"'] \arrow[dr,phantom,"!"] & F(j) \arrow[d,"F(\lambda)"] \\
F_\alpha(j') \arrow[r,"u_\alpha(j')"] & F(j')
\end{tikzcd}
\tag*{$(*)$}
\]
(indexed by $\alpha\in A$). For fixed $j'$ in $J$, the number of these squares is bounded above by $\card\Ob(J/j')\times\card A\leqslant\pi\cdot\pi=\pi$. It follows, since $\mathcal{E}_{j'}$ is large relative to $\pi$, that there exists an object $X(j')$ of $\mathcal{E}_{j'}$, and an arrow $v(j') : F(j')\longrightarrow X(j')$ in such a way that for every $\alpha\in A$, and every $\lambda\in\Ob J/j'$, the square deduced from $(*)$

\msp{253}
by composing with $F(j')\longrightarrow X(j')$ is commutative. We make this choice for every object $j'$ of $J$. Let us then apply F$'$ 1 to $F$ and to the system of arrows $F(j)\xrightarrow{v(j)}X(j)$, we find $G$ and a system of arrows
\[
w(j)\ :\ X(j) \longrightarrow G(j)
\]
in such a way that the composite arrows
\[
\begin{tikzcd}[column sep=small]
w(j)v(j) : & F(j) \arrow[rr] \arrow[dr,dotted] & & G(j) \\
& & X(j) \arrow[ur,dotted] &
\end{tikzcd}
\]
define a homomorphism of sections. I say that this homomorphism has the desired property for F $3_\pi$, i.e. that for every $\alpha\in A$, the system of composites
\[
F_\alpha(j) \longrightarrow F(j) \longrightarrow G(j)
\]
(for $j\in\Ob J$ variable) defines a homomorphism of sections $F_\alpha\longrightarrow G$. Indeed, for this it suffices to consider, for an arrow $\lambda$ of $J$, the
% foot: Version 17 décembre 2025
\printedpage{96}
% head: A. Grothendieck : Dérivateurs, Ch. XVIII, Édition des Universités de Montpellier II et Paris VII
\msp{254}
diagram
\[
\begin{tikzcd}
F_\alpha(j) \arrow[r,"u_\alpha(j)"] \arrow[d,"F_\alpha(\lambda)"'] \arrow[dr,phantom,"!"] & F(j) \arrow[r,"v(j)"] \arrow[d,"F(\lambda)"'] & X(j) \arrow[r,"w(j)"] \arrow[d,phantom,"\text{comm.}"] & G(j) \arrow[d,"G(\lambda)"] \\
F_\alpha(j') \arrow[r,"u_\alpha(j')"] & F(j') \arrow[r,"v(j')"] & X(j') \arrow[r,"w(j')"] & G(j')\ .
\end{tikzcd}
\]
Lemma 4.10.5 being proved, everything thus comes down to proving condition F 1, and its complement F$'$ 1, which will be proved in the same breath.

\textbf{4.10.6.} So we are given the system of objects
\[
X(j)\ \in\ \Ob\mathcal{E}_j \qquad (j\in\Ob J)\ , \tag*{(4.10.6.1)}
\]
and as the case may be, an object $F$ of $\uGamma(\mathcal{E}/J)$, and a system of arrows
\[
F(j) \xrightarrow{u_j} X(j)\ , \tag*{(4.6.10.2)} % [sic] (4.6.10.2) -> (4.10.6.2), numbering of 4.10.6; cited as (4.10.6.2) below
\]
and we propose to construct $G$ in $\uGamma(\mathcal{E}/J)$, and a system of arrows
\[
X(j) \xrightarrow{w(j)} G(j)\ , \tag*{(4.10.6.3)}
\]
in such a way that (if we have the additional datum of the $u(j)$) the composite arrows
\[
F(j) \xrightarrow{u(j)} X(j) \xrightarrow{v(j)} G(j) % [sic] v(j) -> w(j), the arrow X(j) -> G(j) is w(j) in (4.10.6.3)
\]

\msp{255}
are functorial in $j$.

We are going to make a transfinite construction on the ordinals $\alpha$ such that
\[
\card\alpha\ \leqslant\ \pi
\]
of a direct system
\[
(X_\alpha(j))_\alpha
\]
of objects of $\prod_{j\in\Ob J}\mathcal{E}_j$. We are first going to ignore the possible datum (4.10.6.2). We shall put
\[
X_0(j)\ =\ X(j)\ , \tag*{(4.10.6.4)}
\]
and we are going to give the prescription for passing from $X_\alpha$ to $X_{\alpha+1}$. The construction will also include, for every arrow $\lambda : j'\longrightarrow j$ of $J$ and every $\alpha$, the datum of
\[
X_\alpha(j) \xrightarrow{t_\alpha(\lambda)} X_{\alpha+1}(j)\ , \qquad \text{arrow (in $\mathcal{E}$) over $\lambda$,} \tag*{(4.10.6.5)} % [sic] X_\alpha(j) -> X_\alpha(j'), the source of an arrow over \lambda : j' -> j (cf. p. 97: X_\alpha(j') -> X_{\alpha+1}(j))
\]
and these arrows, and the transition arrows $X_\alpha(j)\xrightarrow{\tau_\alpha(j)}X_{\alpha+1}(j)$, must
% foot: Version 17 décembre 2025
\printedpage{97}
% head: A. Grothendieck : Dérivateurs, Ch. XVIII, Édition des Universités de Montpellier II et Paris VII
\msp{256}
satisfy the following three conditions:
\begin{itemize}[leftmargin=3em]
\item[(i)] If $\lambda$ is an identity arrow $\mathrm{id}_J : J\longrightarrow J$, then \fbox{$t_\alpha(\lambda)=\tau_\alpha(j)$.} % [sic] id_J : J -> J -> id_j : j -> j, the identity arrow of an object j of J
\item[(ii)] In the diagram
\[
\begin{tikzcd}
X_\alpha(j') \arrow[r,"\tau_\alpha(j')"] \arrow[dr,"t_\alpha(\lambda)"] & X_{\alpha+1}(j') \arrow[r] \arrow[dr,"t_{\alpha+1}(\lambda)"] \arrow[d,phantom,"\bullet"] & X_{\alpha+2}(j') \arrow[r] \arrow[dr] & X_{\alpha+3}(j') \arrow[r] \arrow[dr,dotted] & {} \\
X_\alpha(j') \arrow[r] & X_{\alpha+1}(j') \arrow[r,"\tau_{\alpha+1}(j')"] & X_{\alpha+2}(j') \arrow[r,"\tau_{\alpha+2}(j')"] & X_{\alpha+3}(j') \arrow[r] & {,}
\end{tikzcd}
\] % [sic] bottom row printed with j' (X_\alpha(j'), X_{\alpha+1}(j'), \tau_{\alpha+1}(j'), X_{\alpha+2}(j'), \tau_{\alpha+2}(j'), X_{\alpha+3}(j')) -> j, the arrows t(\lambda) go to the objects at j (cf. the text below: \tau_{\alpha+1}(j), \tau_{\alpha+2}(j))
$\tau_{\alpha+2}(j)$ ``makes the square $\bullet$ commutative'', i.e. equalizes $t_{\alpha+1}(\lambda)\tau_\alpha(j')$ and $\tau_{\alpha+1}(j)t_\alpha(\lambda)$:
\[
\boxed{\tau_{\alpha+2}(j)t_{\alpha+1}(\lambda)\tau_\alpha(j')\ =\ \tau_{\alpha+2}(j)\tau_{\alpha+1}(j)t_\alpha(\lambda)\ .}
\]
\item[(iii)] If we have two consecutive arrows $\lambda : j''\longrightarrow j'$, $\mu : j'\longrightarrow j$, we thus have the diagram
\[
\begin{tikzcd}
X_\alpha(j'') \arrow[r] \arrow[dr,"t_\alpha(\lambda)"] \arrow[ddr,"t_\alpha(\mu\lambda)"] & X_{\alpha+1}(j'') \arrow[r] \arrow[dr] & X_{\alpha+2}(j'') \arrow[r] \arrow[dr] & X_{\alpha+3}(j'') \\
X_\alpha(j') \arrow[r] & X_{\alpha+1}(j') \arrow[r] \arrow[dr,"t_{\alpha+1}(\mu)"] & X_{\alpha+2}(j') \arrow[r] \arrow[dr] & X_{\alpha+3}(j') \\
X_\alpha(j) \arrow[r] & X_{\alpha+1}(j) \arrow[r,"\tau_{\alpha+1}(j)"] & X_{\alpha+2}(j) \arrow[r,"\tau_{\alpha+2}(j)"] & X_{\alpha+3}(j)\ ,
\end{tikzcd}
\]
\[
\boxed{\tau_{\alpha+2}(j)t_{\alpha+1}(\mu)t_\alpha(\lambda)\ =\ \tau_{\alpha+2}(j)\tau_{\alpha+1}(j)t_\alpha(\mu\lambda)}
\]
\end{itemize}

\msp{257}
Thus, to define the $X_{\alpha+1}(j)$, we have to take account of all these constraints. So:
\begin{itemize}[leftmargin=3em]
\item[1°)] For fixed $j$, $X_{\alpha+1}(j)$ must be the target of the arrows
\[
X_\alpha(j') \xrightarrow{t_\alpha(\lambda)} X_{\alpha+1}(j)\ ,
\]
for the $\lambda\ :\ j'\longrightarrow j$ in $\Ob J/j$ -- including for $\lambda=\mathrm{id}_j$, in which case $t_\alpha(\lambda)\ :\ X_\alpha(j)\longrightarrow X_{\alpha+1}(j)$ will play the role of the transition morphism $\tau_\alpha(j)$.
\item[2°)] If $\alpha$ is of the form $\alpha'+2$, i.e. if it is neither a limit ordinal, nor the successor of a limit ordinal, then we require in addition, for the transition morphisms $\tau_\alpha(j)$, that they equalize the pairs of arrows provided for in (ii) and (iii) above.
\end{itemize}
Since the set of these data in 1°), or of these constraints in 2°),
\msp{258}
has cardinality $\leqslant\pi$, and the category $\mathcal{E}_j$ is large relative to $\pi$, we see that we can indeed find $X_{\alpha+1}(j)$ and $t_\alpha(\lambda)\ :\ X_\alpha(j')\longrightarrow X_{\alpha+1}(j)$, so as to satisfy (if $\alpha=\alpha'{+}2$) the constraints recalled in 2°. This completes the description of the construction of $X_{\alpha+1}(\bullet)$, in terms of $X_\beta(\bullet)$ $(\beta\leqslant\alpha)$ and the $t_\beta(\lambda)$ $(\beta<\alpha)$.

We must finally define $X_\alpha(\bullet)$ for $\alpha$ a limit ordinal. This time, we ask it only to receive the $X_\beta(\bullet)$ for $\beta<\alpha$. Since the set $I_\alpha$ of these $\beta$ has cardinality $\leqslant\pi$, and the $\mathcal{E}_j$ are large relative to $\pi$, we can indeed find the $X_\alpha(j)$.
% foot: Version 17 décembre 2025

\printedpage{98}
% head: A. Grothendieck : Dérivateurs, Ch. XVIII, Édition des Universités de Montpellier II et Paris VII
\msp{259}
We have thus proved the essential part of

\textbf{Lemma 4.10.7.} \textit{Under the conditions of 4.10.3 a) b), and for a given system $X(\bullet)\in\Ob\prod_{j\in\Ob J}\mathcal{E}_j$, we can find a direct system $(X_\alpha(i))_\alpha$ in $\prod\mathcal{E}_j$, indexed by the ordinals $\alpha$ such that} $\card\alpha\leqslant\pi$, \textit{and for every arrow $\lambda : j'\longrightarrow j$ in $J$ arrows $t_\alpha(\lambda) : X_\alpha(j')\longrightarrow X_\alpha(j)$, in such a way that we have the following conditions.} % [sic] (X_\alpha(i))_\alpha -> (X_\alpha(j))_\alpha, the index is j \in Ob J; [sic] X_\alpha(j') -> X_\alpha(j) -> X_\alpha(j') -> X_{\alpha+1}(j), cf. (4.10.6.5) and p. 97
\begin{itemize}[leftmargin=3em]
\item[\textit{1°)}] \textit{$X_0(\bullet)=X(\bullet)$.}
\item[\textit{2°)}] \textit{The conditions (i) (ii) (iii) (p. 256).}
\item[\textit{3°)}] \textit{If we are given $F$ in \textup{Ob}$\,\uGamma(\mathcal{E}/J)$, and arrows $u(j) : F(j)\longrightarrow X(j)=X_0(j)$, we may suppose that for every arrow $\lambda : j'\longrightarrow j$ in $J$, we have}
\[
\begin{tikzcd}
F(j') \arrow[r,"u(j')"] \arrow[d,"F(\lambda)"'] & X_0(j') \arrow[dr,"t_0(\lambda)"] & \\
F(j) \arrow[r,"u(j)"] & X_0(j) \arrow[r,"\tau_0(j)"] & X_1(j)
\end{tikzcd}
\qquad\qquad \text{commutative.}
\tag*{$(*)$}
\]
\end{itemize}

\msp{260}
It only remains to prove this last complement 3°. This is a constraint only on the construction of $X_1(\bullet)$ and of the $t_0(\lambda)$. So suppose that we have constructed $X_1(\bullet)$ without taking this constraint into account. For fixed $j$, we must therefore ``make commutative'' all these squares, i.e. all these double arrows $(t_0(\lambda)u(j'),\tau_0(j)u(j)F(\lambda))$ with target $X_1(j)$. Since the set of these double arrows has cardinality $\card\Ob\mathcal{E}/j\leqslant\pi$, we can find $X_1(j)\longrightarrow X'_1(j)$ in $\mathcal{E}_j$, equalizing all these double arrows, so at the cost of replacing $X_1(j)$ by $X'_1(j)$, we win.

\msp{261}
It is finally time to appeal to hypothesis c) of 4.10.3. So let $\alpha$ be a limit ordinal, as stipulated in this condition. Since $\card\alpha\leqslant\pi$, a fortiori $\card\beta\leqslant\pi$ if $\beta<\alpha$, we thus have a direct system
\[
(X_\beta(\bullet))_{\beta\in I_\alpha}
\]
in $\prod_{j\in J}\mathcal{E}_j$. Since the $\varinjlim$ of type $I_\alpha$ exist in $\mathcal{E}_j$, we can therefore consider
\[
G(j)\ =\ \varinjlim_{\beta\in I_\alpha} X_\beta(j)\ \in\ \Ob\mathcal{E}_j\ .
\]
If $\lambda : j'\longrightarrow j$ is an arrow of $J$, we deduce from it a $\lambda$-arrow
\[
G(\lambda)\ :\ G(j)=\varinjlim X_\beta(j) \longrightarrow G(j')\simeq\varinjlim X_{\beta+1}(j')
\]
by $G(\lambda)=\varinjlim t_\beta(\lambda)$. We must
\msp{262}
take note that the $t_\beta(\lambda)$ do not define a homomorphism of inverse systems, but by condition (ii), p. 256, they define a homomorphism of ind-objects, so
% foot: Version 17 décembre 2025

\end{document}
